\documentclass[a4paper,12pt]{article}

\input{../Preambulo.tex}
\newcommand{\assunto}{FUNÇÕES EXPONENCIAIS}
\newcommand{\atividade}{prova} % prova ou lista
\newcommand{\turma}{ELETROELETRÔNICA 1.\textordmasculine\ ANO}
\newcommand{\numProva}{A}
\newcommand{\professor}{Igor Martins Silva}
\newcommand{\bimestre}{3}
\newcommand{\ano}{2026}
\newcommand{\mesTexto}{agosto}
\newcommand{\mesNum}{08}
\newcommand{\dia}{27}
\newcommand{\dataTexto}{\dia\ de \mesTexto\ de \ano}
\newcommand{\dataNun}{\dia/\mesNum/\ano}

\newcommand{\titulo}{}

\ifdefstring{\atividade}{prova}
  {\renewcommand{\titulo}{Avaliação de Matemática}}
{
\ifdefstring{\atividade}{lista}
  {\renewcommand{\titulo}{Lista de exercícios de Matemática}}
{
\ifdefstring{\atividade}{as}
  {\renewcommand{\titulo}{Avaliação Somativa}}
{
\ifdefstring{\atividade}{ad}
  {\renewcommand{\titulo}{Avaliação Diagnóstica}}
  {\renewcommand{\titulo}{Atividade de Matemática}}
}}}

\newcommand{\emtipo}{%
  \ifdefstring{\tipo}{individual}{\tipo}{em \tipo}%
}

\edef\pathCapa{\currfiledir}

\newcommand{\subtitulobase}[3]{%
    % #1 = título da 1ª coluna
    % #2 = conteúdo da 1ª coluna
    % #3 = mostra número da prova? (vazio ou algo)
    
    \noindent
    \begin{minipage}{2.8cm}
        \includegraphics[width=2.8cm]{\pathCapa Logo_Cefet.png}
    \end{minipage}
    \hfill
    \begin{minipage}{\dimexpr\textwidth-6.2cm\relax}
        \begin{center}
        CEFET-MG -- Campus Contagem \\[3pt]
        \titulo\ -- \bimestre.\textordmasculine\ bimestre de \ano \\[3pt]
        \professor
    \end{center}
    \end{minipage}
    \hfill
    \begin{minipage}{3.2cm}
        \includegraphics[width=3.2cm]{\pathCapa Brasao_Brasil.png}
    \end{minipage}
    
    \begin{center}
        \vspace{15pt}
        \begin{tabular}{|C{210pt}|C{70pt}|C{210pt}|}
            \hline
            \textbf{#1} &
            \textbf{DATA} &
            \textbf{TURMA} \\
            \hline
            #2 & \dataNun & \turma \\
            \hline
        \end{tabular}
    \end{center}
    
    #3
    
    \vspace{15pt}
}

\newcommand{\subtitulolis}{%
  \subtitulobase{ASSUNTO}{\assunto}{}%
}

\newcommand{\subtitulopr}{%
  \subtitulobase{ASSUNTO}{\assunto}{\boxed{\numProva}}%
}

\newcommand{\subtituloas}{%
  \subtitulobase{DISCIPLINA}{MATEMÁTICA}{\boxed{\numProva}}%
}

\newcommand{\subtituload}{%
  \subtitulobase{DISCIPLINA}{MATEMÁTICA}{}%
}

\newcommand{\valorquestao}{}

\newcommand{\definevalorquestao}{%
    \ifdefstring{\atividade}{lista}{%
        % Não faz nada para lista
    }{%
        \ifdefstring{\modo}{objetivo}{%
            \renewcommand{\valorquestao}{\ValorQObj ponto}
        }{%
            \renewcommand{\valorquestao}{\ValorQDisc pontos}
        }%
    }%
}

\newenvironment{exercicioBanco}[1][]{%
    \ifdefstring{\atividade}{lista}{%
        \begin{exercicio}%
    }{%
        \begin{exercicio}[#1]%
    }%
}{%
    \end{exercicio}
}

\ifdefstring{\atividade}{lista}{
    \pagecolor{background-color}
}{
    
}



\hypersetup{
    pdfauthor={\professor},
    pdftitle={Estatística da prova \ -- \assunto \ -- \turma},
}

\title{Estatística da prova \ -- \assunto \ -- \turma}
\author{\professor}
\date{\data}


%************* INÍCIO DO DOCUMENTO *************
\begin{document}

\edef\pathCapa{\currfiledir}
    
    %
    \noindent
    \begin{minipage}{2.8cm}
        \includegraphics[width=2.8cm]{\pathCapa Logo_Cefet.png}
    \end{minipage}
    \hfill
    \begin{minipage}{\dimexpr\textwidth-6.2cm\relax}
    \begin{center}
        CEFET-MG -- Campus Contagem \\[3pt]
        Estatística da prova\ -- \bimestre.\textordmasculine\ Bimestre de \ano \\[3pt]
        \professor
    \end{center}
    \end{minipage}
    \hfill
    \begin{minipage}{3.2cm}
        \includegraphics[width=3.2cm]{\pathCapa Brasao_Brasil.png}
    \end{minipage}
    %
    
    \begin{center}
        \vspace{15pt}
        \begin{tabular}{|C{210pt}|C{70pt}|C{210pt}|}
            \hline
            \textbf{ASSUNTO} &
            \textbf{DATA} &
            \textbf{TURMA} \\
            \hline
            \assunto & \dataNun & \turma \\
            \hline
        \end{tabular}
    \end{center}
    
    \vspace{20pt}
    
    \begin{table}[H]
        \centering
        \renewcommand{\arraystretch}{1.1}
        \caption{Estatística descritiva da prova.}
        \begin{tabular}{C{40pt}C{110pt}C{110pt}C{50pt}C{50pt}C{30pt}C{50pt}}
            \toprule
            Total de alunos & Menor nota (número de alunos com essa nota) & Maior nota (número de alunos com essa nota) & Média da nota & Mediana & Moda & Desvio padrão \\
            \midrule
            32 & 0,6 (1) & 4 (1) & 2,53 & 2,8 & 1,4 & 0,87  \\
            \bottomrule
        \end{tabular}
    \end{table}
    
    \begin{table}[H]
        \centering
        \renewcommand{\arraystretch}{1.1}
        \caption{Frequência de erros em cada questão objetiva.}
        \begin{tabular}{C{50pt}C{80pt}C{80pt}C{80pt}C{80pt}C{80pt}}
            \toprule
            Prova A & Questão 1 & Questão 2 & Questão 3 & Questão 4 & Questão 5 \\
            \midrule
            Prova B & Questão 4 & Questão 5 & Questão 1 & Questão 2 & Questão 3 \\
            \midrule
            & \(\dfrac{4}{32} = 12,5\%\) & \(\dfrac{4}{32} = 12,5\%\) & \(\dfrac{24}{32} = 75\%\) & \(\dfrac{2}{32} \approx 6,3\%\) & \(\dfrac{25}{32} \approx 78,1\%\) \\
            \cmidrule[0.08em](l{2pt}r{0pt}){2-6}
        \end{tabular}
    \end{table}
    
    \begin{figure}[H]
        \centering
        \caption{Histograma e boxplot da nota.}
        \begin{subfigure}[b]{0.5\textwidth}
            \centering
            \includesvg[width=\textwidth]{histograma_Notas}
        \end{subfigure}
        \hfill
        \begin{subfigure}[b]{0.3\textwidth}
            \centering
            \includesvg[width=\textwidth]{boxplot_Notas}
        \end{subfigure}
        \label{H&B_Q1}
    \end{figure}
    
    \begin{flushleft}
        \textbf{INTERPRETAÇÃO}
    \end{flushleft}

    A média de 2,53 corresponde a 63,3\% da pontuação máxima. O coeficiente de variação de Pearson de 0,34 indica dispersão alta das notas, e o coeficiente de assimetria de Pearson de -0,93 indica assimetria moderada à esquerda. O boxplot mostra mínimo de 0,6 (1 aluno) e máximo de 4 (1 aluno), com 50\% das notas concentradas entre 1,7 e 3,2. Complementarmente, o histograma revela que a classe modal é o intervalo \([3,0;3,5)\), concentrando 28,1\% dos alunos, enquanto cerca de 43,7\% ficaram abaixo de 2,5. Quanto às questões objetivas, a questão 5 (prova A) apresentou o maior índice de erro (78,1\%), enquanto a questão 4 (prova A) apresentou o menor (6,3\%).
    
\end{document}
%*************** FINAL DO DOCUMENTO ***************
